\subsection{Two exact limitations of capacity normalization}
\label{subsec:capacity-pullback-limitations}

The critical constants appearing in the reciprocal and Hankel constructions
are stable under two natural operations.  The statements below are deliberately
limited to polynomial pullbacks of the standard adelic set and to changes of
basis of the full Hankel determinant.  They do not constitute a no-go theorem
for weighted, place-dependent, or nonlinear auxiliary constructions.  This is
the boundary at which the adelic-capacity mechanism used in, for example,
Dimitrov's Schinzel--Zassenhaus argument ceases to improve the present
normalization \cite{Dimitrov2019}; we use the standard local capacity
transformation law in the form reviewed by Bost and Chambert-Loir
\cite{BostChambertLoir2009}.

We normalize the absolute values of \(\Q\) by
\(|p|_p=p^{-1}\) and the usual absolute value at infinity.  At the
archimedean place put
\[
 \mathcal E_\infty=[-2,2],
\]
and, at every finite place, let
\[
 \mathcal E_p=
 \bigl\{z\in\mathbb A_{\C_p}^{1,\mathrm{an}}:|z|_p\leq1\bigr\}
\]
be the full closed Berkovich unit disk.  Capacities below are logarithmic
capacities relative to infinity.

\begin{proposition}[Polynomial pullbacks preserve critical adelic capacity]
\label{prop:adelic-polynomial-capacity}
One has
\[
 \capc_\infty(\mathcal E_\infty)=1,
 \qquad
 \capc_p(\mathcal E_p)=1
 \quad\text{for every prime }p.
\]
Consequently the global capacity of the adelic set
\(\mathcal E=(\mathcal E_v)_{v\leq\infty}\) is
\[
 \prod_{v\leq\infty}\capc_v(\mathcal E_v)=1.
\]
More generally, let
\[
 R(X)=a_dX^d+\cdots+a_0\in\Q[X],
 \qquad d\geq1,\quad a_d\ne0.
\]
Then, at every place \(v\),
\begin{equation}\label{eq:local-capacity-pullback}
 \capc_v\bigl(R^{-1}(\mathcal E_v)\bigr)
 =
 \left(
   \frac{\capc_v(\mathcal E_v)}{|a_d|_v}
 \right)^{1/d},
\end{equation}
and hence
\begin{equation}\label{eq:global-capacity-pullback}
 \boxed{
 \prod_{v\leq\infty}
 \capc_v\bigl(R^{-1}(\mathcal E_v)\bigr)=1.}
\end{equation}
Thus rational affine changes, rational polynomial pullbacks, and in
particular rational Chebyshev pullbacks cannot turn this specific critical
adelic capacity into a strict inequality below one.
\end{proposition}

\begin{proof}
The classical identity
\(\capc([-2,2])=1\) gives the archimedean assertion, and the closed
Berkovich disk of radius one has nonarchimedean capacity one.  The local
polynomial transformation law for logarithmic capacity gives
\eqref{eq:local-capacity-pullback}.  Multiplying it over all places yields
\[
 \prod_v\capc_v\bigl(R^{-1}(\mathcal E_v)\bigr)
 =
 \left(
   \frac{\prod_v\capc_v(\mathcal E_v)}
        {\prod_v|a_d|_v}
 \right)^{1/d}.
\]
The first product in the numerator is one, while the product formula for
\(a_d\in\Q^\times\) gives \(\prod_v|a_d|_v=1\).  This proves
\eqref{eq:global-capacity-pullback}.
\end{proof}

\begin{remark}[Why the finite-place set is not \(\mathbb Z_p\)]
The full Berkovich unit disk in Proposition~\ref{prop:adelic-polynomial-capacity}
must not be replaced silently by the classical compact subset
\(\mathbb Z_p\subset\Q_p\).  With the standard logarithmic kernel,
\begin{equation}\label{eq:capacity-Zp}
 \capc_p(\mathbb Z_p)=p^{-1/(p-1)},
\end{equation}
not one.  Indeed, if \(\mu_p\) is normalized Haar measure on
\(\mathbb Z_p\), then
\[
 \int_{\mathbb Z_p^2}v_p(x-y)\,d\mu_p(x)d\mu_p(y)
 =\sum_{k\geq1}\Pr\bigl(v_p(x-y)\geq k\bigr)
 =\sum_{k\geq1}p^{-k}=\frac1{p-1}.
\]
The corresponding equilibrium energy is
\(\log(p)/(p-1)\), which gives \eqref{eq:capacity-Zp}.
In particular, the literal family with finite components \(\mathbb Z_p\)
does not have global capacity one.
\end{remark}

We next make the critical cancellation in
Theorem~\ref{thm:hankel} integral.  For \(n\geq1\), put
\begin{equation}\label{eq:hankel-integral-clearing}
 \mathcal B_n
 =n!\prod_{j=1}^n
   \binom{2j}{j}\binom{2j-1}{j},
 \qquad
 \ell_n=\operatorname{lcm}(1,\ldots,n),
 \qquad
 \mathcal I_n=\ell_n\mathcal B_n.
\end{equation}

\begin{proposition}[Integral Hankel saturation and full-determinant rigidity]
\label{prop:hankel-integral-saturation}
For every \(n\geq1\) and every \(x\in\C\),
\begin{equation}\label{eq:hankel-cleared-linear-form}
 \mathcal I_nD_n(x)
 =(-1)^n\bigl(2\ell_nH_n-\ell_nx\bigr).
\end{equation}
In particular, the two coefficients on the right side are integers.  One
has
\begin{equation}\label{eq:hankel-clearing-asymptotics}
 \mathcal B_n^{1/n^2}\longrightarrow4,
 \qquad
 \mathcal I_n^{1/n^2}\longrightarrow4.
\end{equation}
For every fixed \(x\in\C\),
\begin{equation}\label{eq:hankel-critical-product}
 |D_n(x)|^{1/n^2}\longrightarrow\frac14,
 \qquad
 \boxed{
 |\mathcal I_nD_n(x)|^{1/n^2}\longrightarrow1.}
\end{equation}

Let
\[
 M_n(x)=(H_{i+j}-x)_{0\leq i,j\leq n}.
\]
If \(U_n,V_n\in\operatorname{GL}_{n+1}(\Z)\), then
\begin{equation}\label{eq:unimodular-hankel-invariance}
 \det\bigl(U_nM_n(x)V_n\bigr)=\pm D_n(x).
\end{equation}
Hence unimodular row and column lattice reduction cannot alter the
capacity constant of the full determinant.
\end{proposition}

\begin{proof}
Theorem~\ref{thm:hankel} gives
\[
 D_n(x)=(-1)^n\frac{2H_n-x}{\mathcal B_n}.
\]
Multiplication by \(\mathcal I_n=\ell_n\mathcal B_n\) proves
\eqref{eq:hankel-cleared-linear-form}; the integrality follows from
\(\ell_nH_n\in\Z\).  The Stirling estimate already used in the proof of
Theorem~\ref{thm:hankel} gives
\[
 \log\mathcal B_n=(\log4)n^2+O(n\log n).
\]
The elementary Chebyshev bound \(\log\ell_n=O(n)\) therefore proves
\eqref{eq:hankel-clearing-asymptotics}.  Since
\(H_n=\log n+O(1)\), for fixed \(x\) one has
\[
 \log|2H_n-x|=O(\log\log n)
\]
for all sufficiently large \(n\).  Formula
\eqref{eq:hankel-critical-product} now follows either directly or together
with Theorem~\ref{thm:hankel}.  Finally,
\[
 \det(U_nM_n(x)V_n)=\det(U_n)\det(V_n)D_n(x),
\]
and both integer determinants are \(\pm1\), proving
\eqref{eq:unimodular-hankel-invariance}.
\end{proof}

\begin{remark}[Exact scope of the limitation]
Proposition~\ref{prop:hankel-integral-saturation} concerns the covolume, or
equivalently the full determinant, of this particular Hankel lattice.
It does not exclude a short individual row or linear form produced by
lattice reduction, a useful proper minor, a change of the underlying
lattice, an \(x\)-dependent auxiliary construction, or a nonlinear
combination of the moments.  Likewise,
Proposition~\ref{prop:adelic-polynomial-capacity} does not exclude weighted
capacity, place-dependent changes of the local compact sets, proper local
subsets, rational maps with poles, or nonpolynomial pullbacks.  The proved
obstruction is precisely that rational polynomial pullback and unimodular
full-determinant normalization preserve the existing critical equality;
no broader capacity or lattice no-go is asserted.
\end{remark}
